\section{From Prediction to Deployment Decisions}

The holdout establishes workload screening under one fixed configuration; it
does not yet establish that the sandbox can choose a configuration.  For a
frozen scenario $s$, let $F(s)$ be the configurations that satisfy measured
TTFT, TPOT, throughput, and quality constraints, and let $P(s)$ be their
measured energy--performance Pareto set.  Predictions produce
$\widehat F(s)$ and $\widehat P(s)$.  We require three additional gates before
using them for deployment search.

\textbf{Ranking.}  Spearman correlation measures broad ordering, but a search
must also recover the lowest-energy feasible points.  We will report pairwise
accuracy, top-$k$ recovery, Pareto precision/recall, and normalized hypervolume
error.  Every compared method receives the same measured-trial budget and the
same frozen candidate set.

\textbf{SLO safety.}  A false-safe point is predicted to satisfy every SLO but
fails at least one target measurement.  We will report both its candidate rate
and whether the final recommendation is false-safe.  Prediction intervals must
be calibrated on a validation split and frozen again before a disjoint
holdout; the current \DiagnosticCoverage{} diagnostic containment is not used
as a calibrated guarantee.

\textbf{Net measurement cost.}  For measured record $e$ using $g_e$ GPUs for
$t_e$ seconds, purchased cost is
\begin{equation}
C(D)=\sum_{e\in D}g_e t_e/3600.
\end{equation}
The ledger includes calibration, warmup, failed jobs, and final verification.
A sandbox succeeds only if it matches target-only search at the same SLO safety
while buying fewer target GPU-hours; simulator queries alone are not counted as
saved energy.  Table~\ref{tab:roadmap} keeps those claims separate.

\begin{table}[t]
\centering
\caption{Claim ladder.  Completed workload evidence cannot be reused as proof
of configuration, topology, or hardware transfer.}
\label{tab:roadmap}
\scriptsize
\setlength{\tabcolsep}{2.6pt}
\begin{tabular}{@{}p{0.22\columnwidth}p{0.31\columnwidth}p{0.38\columnwidth}@{}}
\toprule
Claim & Real evidence & Gate \\
\midrule
Measurement path & 3-run anchor + raw traces & Boundaries + hashes: done \\
Workload transfer & $6\!\times\!3$ dev + $(8+9)\!\times\!3$ blind & Energy replicated; TTFT scoped \\
Serving controls & Batch/cache grid & SLO safety, Pareto recall \\
TP/PP topology & Sparse scale probes + reruns & Rank and false-safe rate \\
Hardware transfer & H100/H200/B200 anchors & Cold-start cost and target error \\
\bottomrule
\end{tabular}
\end{table}

